\documentclass{aastex701}

% DRAFT 2026-10-13 -- NOT SUBMITTED. Ricardo must read and check every number, the figure
% and the references himself before submission. Do not add a statement that the author
% verified the content unless that is true at the time of submission.

\shorttitle{A candidate 18.9-day planet around TOI-4342}
\shortauthors{Pinho}

\begin{document}

\title{A Candidate Third Transiting Planet on an 18.9-day Orbit around the M Dwarf TOI-4342, Matching a Dynamical Prediction}

\author[orcid=0009-0006-2744-2599,gname=Ricardo,sname=Pinho]{Ricardo Pinho}
\affiliation{Independent researcher}
\email{pinho.ricardo@gmail.com}

\begin{abstract}
I report a candidate third transiting planet around TOI-4342 (TIC 354944123), an M0 dwarf hosting two sub-Neptunes at 5.54 and 10.69 days. Nine sectors of TESS 2-minute photometry show a $1182\pm107$\,ppm, 2.4\,hr transit every 18.908226\,days, $R_p=2.25\pm0.13\,R_\oplus$, in all ten observed windows (signal-to-noise ratio 11.7). Period, radius and inclination match the planet predicted for this system by \citet{2025AJ....170....3T}. Difference images place the source within $11''$ of the target and exclude every Gaia neighbor able to mimic it. With a speckle contrast curve and the pixel-excluded neighbors removed, TRICERATOPS gives a false-positive probability of 0.018, just above the 0.015 validation threshold, and a nearby-false-positive probability of zero. The three planets are coplanar to within $1^\circ$ and dynamically stable. I provide an ephemeris for ground-based confirmation.
\end{abstract}

\keywords{\uat{Exoplanet detection methods}{489} --- \uat{Exoplanets}{498} --- \uat{M dwarf stars}{982} --- \uat{Transit photometry}{1709}}

\section{Introduction}

TOI-4342 ($T=11.03$, 61.5\,pc; $T_{\rm eff}=3866\pm92$\,K, $R_\star=0.598\pm0.018\,R_\odot$, $M_\star=0.587\pm0.013\,M_\odot$) hosts planets b ($P=5.538$\,d) and c (10.689\,d) near the 2:1 resonance \citep{2023AJ....165...93T}, with masses $7.3\pm1.3$ and $4.8\pm1.4\,M_\oplus$ \citep{2026A&A...708A..81P}. \citet{2025AJ....170....3T} predicted an additional planet at $18.5^{+2.7}_{-2.0}$\,d with $2.24\,R_\oplus$. TOI-4342 entered my search for additional transits around small-star planet hosts.

\section{Data and detection}

TESS \citep{2015JATIS...1a4003R} observed TOI-4342 in Sectors 13, 27, 66, 67, 93, 94 and 101--103 (2019--2026). I used SPOC 2-minute PDCSAP light curves and target pixel files \citep{2016SPIE.9913E..3EJ}, 20-second light curves (Sectors 93--103), and QLP \citep{2020RNAAS...4..204H} and TESS-SPOC \citep{2020RNAAS...4..201C} full-frame-image light curves. Each sector was detrended with a 0.75-day biweight filter \citep{2019AJ....158..143H} with all transits masked, and models of b and c were subtracted. A box least-squares search \citep{2002A&A...391..369K} found the 18.9-day signal.

Fourteen predicted transits fall within the sectors; ten have data, including one in cadences that SPOC excludes from its planet search and one partial transit. All ten show the transit (Figure~\ref{fig:1}b): depths agree ($\chi^2=8.7$ for 9 degrees of freedom), as do the 2019--2023 and 2025--2026 halves ($0.98\pm0.11$ and $0.89\pm0.12$ of the model depth). The signal is present in simple-aperture flux ($10.7\sigma$) and in QLP ($8.8\sigma$). Applied to the same data with the real windows removed, the identical statistic never reached the observed value at 3,000 random periods (maximum 8.9 against 13.4). Depths at $P/2$, $P/3$ and $P/4$ are consistent with zero ($53\pm78$\,ppm at $P/2$), and odd and even transits agree ($0.5\sigma$), which excludes $2P$. Transit times are consistent with a linear ephemeris ($\chi^2=7.9$, 6 degrees of freedom).

A {\tt batman} \citep{2015PASP..127.1161K} and {\tt emcee} \citep{2013PASP..125..306F} fit with a stellar-density prior of $3.86\pm0.36$\,g\,cm$^{-3}$ gives $P=18.908226\pm0.000042$\,days, $T_0=2459904.5666\pm0.0021$ BJD$_{\rm TDB}$, $R_p/R_\star=0.0344\pm0.0016$, $b=0.76\pm0.03$, $i=88.96\pm0.07^\circ$, $a/R_\star=41.7\pm1.3$ and a duration of $2.44\pm0.11$\,hr (Figure~\ref{fig:1}a). Fitted freely, the implied density, $8.8^{+4.7}_{-5.9}$\,g\,cm$^{-3}$, is consistent with the star. The inclination matches b ($88.83\pm0.22^\circ$) and c ($89.61\pm0.26^\circ$).

\section{Vetting}

Gaia DR3 \citep{2016A&A...595A...1G,2023A&A...674A...1G} shows no sign of a companion (RUWE 0.98). Gemini-South/Zorro speckle imaging (\citealt{2023AJ....165...93T}; instrument: \citealt{Scott2021}) detected none at $\Delta m_{832}>4.7$ beyond $0.1''$. On a 1976 photographic plate, when the star was $7.4''$ from its 2026 position, no source appears at that position to $\Delta T\approx7$; a 1182\,ppm signal requires $\Delta T<7.4$.

For each sector I fitted the in-transit difference image with that sector's out-of-transit image, shifted across a sky grid, and summed $\chi^2$ over sectors. On the same pixels, b and c give target $\Delta\chi^2$ of 9.5 and 3.5. The candidate gives 11.4, a 3$\sigma$ region within $11''$ of the target, and $\Delta\chi^2\ge79$ for every Gaia star that could mimic it (Figure~\ref{fig:1}c). Its depth versus aperture size tracks those of b and c.

\section{Statistical validation and dynamics}

TRICERATOPS \citep{2021AJ....161...24G}, with SPOC apertures from nine sectors and the Zorro contrast curve, gives a false-positive probability (FPP) of $0.0161\pm0.0002$ and a nearby-false-positive probability (NFPP) of 0.0005. Removing the pixel-excluded neighbors gives FPP $=$ $0.018\pm0.002$ and NFPP $=0$. Adding MOLUSC \citep{2021AJ....162..128W} limits on bound companions from the imaging and Gaia gives $0.016\pm0.011$. The remainder is almost entirely a planet transiting an unresolved bound companion. The candidate therefore sits at, not below, the validation threshold of FPP $<0.015$ and NFPP $<0.001$ \citep{2021AJ....161...24G}, and I do not claim validation. This calculation does not include the lower false-positive rate of candidates in systems that already host transiting planets, the ``multiplicity boost'' used to validate Kepler multis \citep{2012ApJ...750..112L,2014ApJ...784...44L}; that factor would lower the FPP further, but it was calibrated on Kepler's population, and I do not apply it here.

N-body models predict transit-timing variations of a few minutes, below the TESS timing precision of $\approx$8\,minutes per transit; the observed times neither require nor exclude them. SPOCK \citep{2020PNAS..11718194T} with REBOUND \citep{2012A&A...537A.128R} gives a stability probability of 0.89 (median over mass and eccentricity draws; 97\% of draws above 0.5) for b, c and the candidate, including the 47.5-day radial-velocity candidate of \citet{2026A&A...708A..81P}.

\section{A second, weaker signal}

The same data show a $287\pm40$\,ppm transit every 1.957690\,days ($1.11\,R_\oplus$, signal-to-noise ratio 8.8), consistent across seasons, with no odd--even difference and with the expected density. Its pixel localization is looser and TRICERATOPS gives FPP $=$ 0.013 (0.007 with MOLUSC) and NFPP $=0.001$, but its signal-to-noise ratio is close to the level at which my search expects noise detections. I regard it as a lower-confidence candidate.

\section{Follow-up}

TESS will not observe TOI-4342 again through Sector 134. A 1.2\,ppt, 2.4\,hr transit is within reach of 1-m-class telescopes. Predicted mid-transit times are $T_n=2459904.5666+18.908226\,n$ BJD$_{\rm TDB}$, with uncertainty $\approx$12\,minutes through 2027. Transits fully observable from the south include 2026 October 28 02:10 (Chile), 2027 May 24 01:56 (South Africa) and 2027 August 26 14:55\,UTC (Australia).

\section{Data and code}

The TESS data are available at MAST: 2-minute light curves \dataset[doi:10.17909/t9-nmc8-f686]{https://doi.org/10.17909/t9-nmc8-f686} and TESS-SPOC FFI light curves \dataset[doi:10.17909/t9-wpz1-8s54]{https://doi.org/10.17909/t9-wpz1-8s54}.
Code to reproduce this analysis is available at \url{https://github.com/rpinho/planet-hunt} (branch \texttt{round6-toi4342}).

\begin{figure}[ht!]
\plotone{fig1.pdf}
\caption{(a) SPOC 2-minute photometry of nine sectors, with b, c and the 1.96-day signal removed, folded on the 18.9-day ephemeris (15-minute bins), with the best-fit model. (b) Depth at each predicted transit in 2-minute (filled) and QLP (open, offset by 25 days) photometry; the line is the model depth. (c) Difference-image localization summed over nine sectors, $\sqrt{\Delta\chi^2}$ relative to the best position; star: TOI-4342; circles: Gaia stars bright enough to mimic the signal.\label{fig:1}}
\end{figure}

\begin{acknowledgments}
This paper includes data collected by the TESS mission. Funding for the TESS mission is provided by the NASA's Science Mission Directorate. This research has made use of the Exoplanet Follow-up Observation Program (ExoFOP; DOI: 10.26134/ExoFOP5) website, which is operated by the California Institute of Technology, under contract with the National Aeronautics and Space Administration under the Exoplanet Exploration Program. % The Zorro contrast curve is published in Tey et al. (2023), Fig. 5 (program GS-2021A-LP-105); no observer permission needed.
Some of the observations in this paper made use of the High-Resolution Imaging instrument Zorro at Gemini-South (program GS-2021A-LP-105; published by \citealt{2023AJ....165...93T}), retrieved from ExoFOP. Zorro was funded by the NASA Exoplanet Exploration Program and built at the NASA Ames Research Center by Steve B. Howell, Nic Scott, Elliott P. Horch, and Emmett Quigley. This work has made use of data from the European Space Agency (ESA) mission {\it Gaia} (\url{https://www.cosmos.esa.int/gaia}), processed by the {\it Gaia} Data Processing and Analysis Consortium (DPAC, \url{https://www.cosmos.esa.int/web/gaia/dpac/consortium}). Funding for the DPAC has been provided by national institutions, in particular the institutions participating in the {\it Gaia} Multilateral Agreement.
\end{acknowledgments}

\section*{AI Disclosure Statement}
Generative AI was used extensively. Claude Opus 5.5 (Anthropic) wrote and ran the analysis code, carried out the data retrieval, vetting, statistical and dynamical tests, searched the literature, and drafted the text, figure and bibliography.

\facilities{TESS, Gaia, Gemini:South (Zorro), ExoFOP}

\software{astropy \citep{2022ApJ...935..167A}, batman \citep{2015PASP..127.1161K}, emcee \citep{2013PASP..125..306F}, lightkurve \citep{2018ascl.soft12013L}, matplotlib \citep{2007CSE.....9...90H}, numpy \citep{2020Natur.585..357H}, REBOUND \citep{2012A&A...537A.128R}, scipy \citep{2020NaMet..17..261V}, SPOCK \citep{2020PNAS..11718194T}, MOLUSC \citep{2021AJ....162..128W}, tess-point \citep{2020ascl.soft03001B}, TRICERATOPS \citep{2021AJ....161...24G}, wotan \citep{2019AJ....158..143H}}

\bibliography{rnaas_toi4342}{}
\bibliographystyle{aasjournalv7}

\end{document}
